\begingroup
\renewcommand{\arraystretch}{1.45}
\fancyhead[L]{\small ARITHMETIC–GEOMETRIC MEAN}
\Needspace{39\baselineskip}
\AGMArticle{7}

On setting
\[
\frac1{\mathrm{M}(1+x,1-x)}=1+Axx+Bx^4+Cx^6+\text{etc.}
\]
and, as in the preceding article, \(x=\frac{2t}{1+tt}\), we have
\[
\begin{aligned}
&1+A(\frac{2t}{1+tt})^2+B(\frac{2t}{1+tt})^4\\&\quad+C(\frac{2t}{1+tt})^6+D(\frac{2t}{1+tt})^8+\text{etc.}\\[6pt]
&=1+tt+At^4+At^6+Bt^8+Bt^{10}\\&\quad+Ct^{12}+Ct^{14}+Dt^{16}+Dt^{18}+\text{etc.}
\end{aligned}
\]
whence the equations
\[
\begin{aligned}
4A&=1\\[6pt]
-8A+16B&=A\\[6pt]
12A-64B+64C&=A\\[6pt]
-16A+160B-384C+256D&=B\qquad\text{etc.}
\end{aligned}
\]
\Needspace{31\baselineskip}
and hence \(A=\frac14,\allowbreak \ B=\frac14.\frac9{16},\allowbreak \ C=\frac14.\frac9{16}.\frac{25}{36},\allowbreak \ D=\frac14.\frac9{16}.\frac{25}{36}.\frac{49}{64}\), as above;
but the basis of the law would be more laborious to deduce from this, so we prefer the following method. From the equation
\[
\begin{aligned}&\frac{2t}{1+tt}+A(\frac{2t}{1+tt})^3+B(\frac{2t}{1+tt})^5+\text{etc.}\\&=2t(1+At^4+Bt^8.\;.\;.)\end{aligned}
\]
the following equations result:
\[
\begin{array}{@{}rcl@{\quad}r@{}}
1&=&1&[1]\\[6pt]
0&=&1-4A&[2]\\[6pt]
A.&=&1-12A+16B&[3]\\[6pt]
0&=&1-24A+80B-64C&[4]\\[6pt]
B&=&1-40A+240B-448C+256D&[5]\\[6pt]
0&=&1-60A+560B-1792C+2304D-1024E&[6]
\end{array}
\]
\SourceLeaf{377}
\Needspace{40\baselineskip}
whose coefficients are readily brought under a general formula: namely, the \(n^{\mathrm{th}}\) equation will be
\[
\begin{aligned}
M={}&1-4A\times\frac{n.n-1}{1.2}\\&\quad+16B\times\frac{n+1.n.n-1.n-2}{1.2.3.4}\\&\quad-64C\times\frac{n+2.n+1.n.n-1.n-2.n-3}{1.2.3.4.5.6}\\[6pt]
&+256D\times\frac{n+3.n+2.n+1.n.n-1.n-2.n-3.n-4}{1.2.3.4.5.6.7.8}-\text{etc.}
\end{aligned}
\]
where \(M\) will be either \(=0\) (when \(n\) is even), or equal to the \(\frac12(n+1)^{\mathrm{th}}\) term of the series
\(1,\ A,\ B,\ C,\ D\) etc. (when \(n\) is odd). From these equations we now deduce the following new ones\HistoricalNote{27}.
\[
\begin{aligned}
\relax [2],\quad&0=1-4A\\[6pt]
4[3]-[1],\quad&4A-1=3-48A+64B\\[6pt]
9[4]-4[2],\quad&0=5-200A+720B-576C\\[6pt]
16[5]-9[3],\quad&16B-9A=7-532A+3696B-7168C\\&\quad+4096D\\[6pt]
25[6]-16[4],\quad&0=9-1116A+12720B\\&\quad-43776C+57600D-25600E
\end{aligned}
\]
\Needspace{41\baselineskip}
etc., where the coefficients are readily brought under a general law. Namely, the \(n^{\mathrm{th}}\) equation will be
\[
\begin{aligned}
&nnN-(n-1)^2L\\[6pt]&=(2n-1)(1-4A\times\frac{3nn-3n+2}{1.2}
\\&\quad+16B\times\frac{n.n-1.5nn-5n+6}{1.2.3.4}\\[6pt]
&-16C\times\frac{n+1.n.n-1.n-2.7nn-7n+12}{1.2.3.4.5.6}\\[6pt]
&+64D\times\frac{n+2.n+1.n.n-1.n-2.n-3.9nn-9n+20}{1.2.3.4.5.6.7.8}-\text{etc.})\quad\text{\textsuperscript{**)}}
\end{aligned}
\]\begingroup\renewcommand{\thefootnote}{**)}\footnotetext{\csname HistoricalNoteText28\endcsname}\endgroup
where the progression of the factors is evident (that is, besides the simple factors, each coefficient contains a quadratic factor of the form\hfill\\{}
\(knn-\nobreak{}kn+\nobreak{}\frac14(kk-\nobreak{}1)\)). These equations
are expressed more simply in the following way, by dividing each right-hand side into two
parts (except in the first equation, which is retained unchanged):
\[
\begin{aligned}
0&=1-4A\\[6pt]
4A-1&=\{3-36A\\[6pt]
&\phantom{=\{}-12A+64B\}\\[6pt]
0&=\{5-180A+400B\\[6pt]
&\phantom{=\{}-20A+320B-576C\}\\[6pt]
16B-9A&=\{7-504A+2800B-3136C\\[6pt]
&\phantom{=\{}-28A+896B-4032C+4096D\}\\[6pt]
0&=\{9-1080A+10800B-28224C+20736D\\[6pt]
&\phantom{=\{}-36A+1920B-15552C+36864D-25600E\}
\end{aligned}
\]
\SourceLeaf{378}
\Needspace{39\baselineskip}
namely, in general, the \(n^{\mathrm{th}}\) equation
\[
\begin{aligned}
&nnN-(n-1)^2L\\[6pt]
&=(2n-1)\Biggl\{\begin{aligned}
1&-3.4A\frac{n. n-1}{1.2}
\\&\quad+5.16B\frac{n+1. n. n-1. n-2}{1.2.3.4}
\\&\quad-7.64C\frac{n+2. n+1.\;.\;.\;. n-3}{1.2.\;.\;.6}\\[6pt]
&-4A\frac{2}{1.2}
\\&\quad+16B\frac{n. n-1.16}{1.2.3.4}
\\&\quad-64C\frac{n+1. n. n-1. n-2.54}{1.2.3.4.5.6}
\end{aligned}\\[6pt]
&\qquad\begin{aligned}
&.\;.\;.\;.\;.\pm k.2^{k-1}K
\frac{n+\frac{k-3}{2}. n+\frac{k-5}{2}. n+\frac{k-7}{2}.\;.\;. n-\frac{k-1}{2}}
{1.2.3.\;.\;.\;. k-1}.\;.\;.\\[6pt]
&.\;.\;.\;.\;.\pm 2^{k-1}K
\frac{n+\frac{k-5}{2}. n+\frac{k-7}{2}.\;.\;.\;. n-\frac{k-3}{2}.\frac14(k-1)^3}
{1.2.\;.\;.\;.\;.\;.\;.\;.\;. k-1}.\;.\;.
\end{aligned}\Biggr\}
\end{aligned}
\]
where \(k\) denotes any odd number without restriction, so that the equations take the form
\[
\begin{aligned}
0&=1-4A\\[6pt]
4A-1&=3(1-4A)-4(9A-16B)\\[6pt]
0&=5(1-4A)-20(9A-16B)+16(25B-36C)\\[6pt]
16B-9A&=7(1-4A)-56(9A-16B)+112(25B-36C)\\&\quad-64(49C-64D)\\[6pt]
0&=9(1-4A)-120(9A-16B)\\&\quad+432(25B-36C)-576(49C-64D)\\[6pt]
&\qquad+256(81D-100E)
\end{aligned}
\]
\Needspace{38\baselineskip}
etc., and in general the \(n^{\mathrm{th}}\) equation
\[
\begin{aligned}
&nnN-(n-1)^2L\\[6pt]
&=(2n-1)(1-4A)\\&\quad-4\frac{2n-1. n. n-1}{1.2.3}(9A-16B)\\&\quad+16\frac{2n-1. n+1. n. n-1. n-2}{1.2.3.4.5}(25B-36C)\\[6pt]
&\qquad-64\frac{2n-1. n+2. n+1. n. n-1. n-2. n-3}{1.2.3.4.5.6.7}(49C-64D)+\text{etc.}
\end{aligned}
\]
where the universality of the evident law follows directly from the calculation. From this, however, it is clear
that necessarily
\[
\begin{gathered}0=1-4A,\quad 0=9A-16B,\quad 0=25B-36C,\\[6pt]
0=49C-64D,\quad 0=81D-100E\end{gathered}
\]
\textit{etc. to infinity,} and thus
\[
\begin{gathered}A=\tfrac14,\\ B=\tfrac14.\tfrac9{16},\\[6pt]
C=\tfrac14.\tfrac9{16}.\tfrac{25}{36},\\[6pt]
D=\tfrac14.\tfrac9{16}.\tfrac{25}{36}.\tfrac{49}{64},\\[6pt]
E=\tfrac14.\tfrac9{16}.\tfrac{25}{36}.\tfrac{49}{64}.\tfrac{81}{100}\end{gathered}
\]
\textit{and so on to infinity.} Q.~E.~D.
\SourceLeaf{379}
\clearpage
\endgroup
